% =====================================================================
\chapter{Conclusion and Future Work}
\label{ch:conclusion}

\section{Conclusion}
This thesis set out to answer a practical question: can a university in
a developing region give its students live, trustworthy, in-their-own-
language information about their bus, at a cost the institution can
actually sustain? The answer, demonstrated on a real fleet rather than
in simulation, is yes.

We designed, implemented, and deployed \system{}, a complete real-time
bus tracking platform comprising a bilingual passenger mobile
application, a web administration console, and a central application
server backed by a 31-model relational schema---all self-hosted on a
single low-cost server. Its distinguishing technical ideas are:

\begin{itemize}
  \item \textbf{Resilient dual-source acquisition.} A fixed,
  vehicle-grade GPS tracker provides fully automatic, driver-independent
  positioning as the primary source, while the driver's smartphone
  serves only as an optional fallback, selected per fix by health-based
  arbitration. In the pilot this was no theoretical safeguard: 64\% of
  trips ran on hardware and 36\% on the fallback, sustaining coverage
  while most of the fleet awaited tracker installation.
  \item \textbf{A training-free geometric ETA estimator.} Projecting the
  live position onto the route polyline and dividing remaining distance
  by a rolling-average speed yields a per-stop arrival estimate with a
  confidence label from the first day of deployment, with every
  prediction logged for future learning-based refinement.
  \item \textbf{An edge-condition-aware passenger experience.} Pre-trip
  vehicle awareness, explicit staleness signalling, offline retention,
  and proximity-based drop-off alerts keep the interface honest and
  informative precisely where naive trackers show a blank or frozen
  map---in both English and Bangla.
\end{itemize}

The seven-week pilot on the operational fleet (9 routes, 3 buses, 38
riders, 71 trips, 46{,}548 GPS reports) established feasibility with a
median position-acquisition latency of 5.9\,s and an 80\% read rate on
delivered notifications. Just as importantly, the pilot taught a lesson
that no simulation would have: roughly 98\% of raw device reports were
redundant idle-state transmissions. The resulting two-sided
optimization---asymmetric device reporting (5\,s moving / 300\,s
parked) and trip-aware server polling (8\,s active / 60\,s idle)---cut
idle-period reporting by up to $60\times$ while leaving live tracking
untouched, and stands as a concrete, transferable design guideline for
any scheduled-fleet tracking deployment: \emph{couple reporting cadence
to vehicle state and timetable, at both the device and the server.}

\section{Future Work}
\label{sec:concl-future}
The deployed system and its accumulating operational data open several
concrete directions:

\begin{enumerate}
  \item \textbf{Learning-based ETA prediction.} Every prediction and
  every arrival is already logged. Once the trajectory history is large
  enough, a data-driven model (Section~\ref{sec:lit-eta}) can be trained
  and evaluated directly against the geometric baseline on the same
  logged data.
  \item \textbf{Fleet-wide tracker installation.} Fitting the remaining
  buses reduces the smartphone path to a rarely needed safety net and
  enables uniform fleet analytics.
  \item \textbf{Ignition-coupled reporting.} Wiring the tracker's ACC
  line would let the device sleep with the engine, completing the
  efficiency work of Section~\ref{sec:eval-efficiency} with
  engine-state-triggered reporting and distance-based thresholds.
  \item \textbf{Full timetable integration.} Associating every trip with
  its scheduled departure enables systematic schedule-adherence
  analytics for the transport office---the administrative counterpart of
  the passenger-facing ETAs.
  \item \textbf{Crowd-sourced occupancy.} At larger scale, passenger
  participation could add ``how full is the bus'' to ``where is the
  bus,'' complementing rather than carrying the positioning system.
  \item \textbf{A formal usability study.} A wider, longer deployment
  with a structured passenger study would quantify the experience
  benefits that the pilot's 80\% notification read rate only hints at.
\end{enumerate}

\section{Closing Remark}
The deepest lesson of this project is methodological. The system's most
valuable finding---the 98\% idle-reporting waste---was not visible in
any design document, simulator, or code review. It appeared only in the
operational data of a real deployment, and once visible, it was fixable
within days. Building the audit layer first, deploying early, and
letting production data drive the engineering agenda turned a student
project into an evaluated system; we commend the same path to those who
follow.
